% !PRINTED 764
\weberpage{764}
4. \gesperrt{The proof of Theorem I. is hereby reduced to the proof of the following special case: If there subsists an equation}

\weq{11}{C_1e^{z_1}+C_2e^{z_2}+\cdots+C_me^{z_m} = 0,}

\gesperrt{in which the $z_1, z_2, \ldots, z_m$ are algebraic numbers different from one another, in which $C_1, C_2, \ldots, C_m$ and the combinations}

\weq{12}{C_1g(z_1)+C_2g(z_2)+\cdots+C_mg(z_m)}

\gesperrt{are rational numbers, when $g(z)$ is any integral function with rational coefficients, then the $C_1, C_2, \ldots, C_m$ must all vanish.}

In order to prove this theorem, we first determine a positive integral rational number $c$ in such a way that

\weq{13}{x_1 = cz_1,\ x_2 = cz_2,\ \ldots,\ x_m = cz_m}

become integral algebraic numbers (§.\,133, 5.). We assume the coefficients $C_1, C_2, \ldots, C_m$ to be \gesperrt{integral} rational numbers, and form an integral function

\weq{14}{\varphi(x) = ax^r+a_1x^{r-1}+\cdots+a_r}

with integral rational numerical coefficients, which has the following properties:

1) The numbers $x_1, x_2, \ldots, x_m$ occur among the roots of $\varphi(x)=0$.

2) The sum

\[
C_1\varphi'(x_1)+C_2\varphi'(x_2)+\cdots+C_m\varphi'(x_m) = k,
\]

which by the assumptions (12), (13) is an integral rational number, is different from zero.

In order to see that such a function $\varphi(x)$ always exists, let one first take a function $\chi(x)$ which satisfies only the condition 1), and the second condition that none of the numbers $x_1, x_2, \ldots, x_m$ is a double root of $\chi(x)$; such a function evidently exists [in order to form a function $\chi$ one may, for example, free the norm of the product $(x-x_1)(x-x_2)\cdots(x-x_m)$ from common divisors with its derivatives]. If one then sets

\[
\varphi(x) = x^h\chi(x),\ \varphi'(x_1) = x_1^h\chi'(x_1),\ \varphi'(x_2) = x_2^h\chi'(x_2)\ldots,
\]
%CONT>>
